[
    node distance=0.8cm and 1.2cm,
    auto,
    block/.style={
        rectangle,
        draw,
        fill=blue!5,
        text width=2.2cm,
        align=center,
        rounded corners,
        minimum height=1.0cm,
        font=\small
    },
    head/.style={
        rectangle,
        draw,
        fill=yellow!10,
        text width=1.8cm,
        align=center,
        rounded corners,
        minimum height=0.6cm,
        font=\footnotesize
    },
    line/.style={
        draw,
        -{Latex[length=2.5mm]},
        thick,
        gray!80
    },
    barrier/.style={
        circle,
        draw=red!80,
        fill=red!10,
        very thick,
        minimum size=0.6cm,
        inner sep=0pt
    },
    success/.style={
        circle,
        draw=green!60!black,
        fill=green!10,
        very thick,
        minimum size=0.6cm,
        inner sep=0pt
    },
    labelnode/.style={
        font=\bfseries\small,
        anchor=west
    }
]

    \node[block] (input1) {Input $P_t$\\ \textit{(Global)}};
    \node[block, right=of input1] (net1) {Single-Head\\ Network};
    \node[barrier, right=of net1] (out1) {\Large $\times$};
    \node[right=0.1cm of out1, align=left, font=\footnotesize, text width=2.5cm] {Cannot compute\\ global parity\\ under FSQI/$\AC^0$};

    \path[line] (input1) -- (net1);
    \path[line] (net1) -- (out1);
    \node[above=0.45cm of input1, labelnode] {(a) Barrier: Single-Frame, Single-Head};

    \node[block, below=1.3cm of input1] (input2) {Input $(P_{t-1}, P_t)$\\ \textit{(Transition)}};
    \node[block, right=of input2] (net2) {Single-Head\\ Network};
    \node[success, right=of net2] (out2) {\Large $\checkmark$};
    \node[right=0.1cm of out2, align=left, font=\footnotesize, text width=2.5cm] {Computes Local\\ Difference $\Delta$};

    \path[line] (input2) -- (net2);
    \path[line] (net2) -- (out2);
    \node[above=0.45cm of input2, labelnode] {(b) Bypass 1: Multi-Frame Input};

    \node[block, below=1.3cm of input2] (input3) {Input $P_t$\\ \textit{(Global)}};

    \node[right=of input3, text width=0cm, inner sep=0] (split) {};
    \node[head, right=0.2cm of split, yshift=0.4cm] (head1) {Head $\pi_0$};
    \node[head, right=0.2cm of split, yshift=-0.4cm] (head2) {Head $\pi_{B-1}$};
    \node[font=\footnotesize] at ($(head1)!0.5!(head2)$) {$\vdots$};

    \node[success, right=0.8cm of head1, yshift=-0.4cm] (out3) {\Large $\checkmark$};
    \node[right=0.1cm of out3, align=left, font=\footnotesize, text width=2.5cm] {Rollout outcomes\\ recover $B$ bits};

    \path[line] (input3) -- (split);
    \draw[line] (split.center) |- (head1.west);
    \draw[line] (split.center) |- (head2.west);
    \draw[line] (head1.east) -| (out3.north);
    \draw[line] (head2.east) -| (out3.south);

    \node[above=0.45cm of input3, labelnode] {(c) Bypass 2: Multiple Rollout Heads};

    \node[block, below=1.3cm of input3] (input4) {Rows of $P_t$\\ then $P_{t+1}$};
    \node[block, right=of input4] (net4) {Recurrent State\\ $s\leftarrow s\oplus x_i$};
    \node[success, right=of net4] (out4) {\Large $\checkmark$};
    \node[right=0.1cm of out4, align=left, font=\footnotesize, text width=2.5cm] {Nim-sum after\\ $N$ row updates\\ (growing latency)};
    \path[line] (input4) -- (net4);
    \path[line] (net4) -- (out4);
    \node[above=0.45cm of input4, labelnode] {(d) Bypass 3: Persistent Recurrence};
